\documentclass[11pt]{article}
\usepackage[margin=1in]{geometry}
\usepackage{amsmath,amssymb}
\usepackage{enumitem}
\usepackage{booktabs}
\usepackage{framed}
\usepackage{bbm}

\newcommand{\indep}{{\bot\negthickspace\negthickspace\bot}}
\DeclareMathOperator{\E}{\mathbb{E}}

\pagestyle{plain}

\begin{document}

\begin{center}
{\Large \textbf{PS200C: Quantitative Methods}} \\[6pt]
{\large In-class Activity 9: Regression Discontinuity Designs}
\end{center}

\vspace{4pt}
\noindent\textbf{Name:} \hrulefill

\vspace{12pt}

%%% SETUP for Q1 %%%
\begin{framed}
\noindent\textbf{Setup for Q1.}  In the United States, Medicare provides nearly universal health insurance to people aged 65 and older.  At age 65 a person becomes eligible for Medicare regardless of their previous insurance status.  A researcher wants to estimate the effect of Medicare on hospital utilization, using age as the running variable.
\end{framed}

\vspace{6pt}

%%% QUESTION 1 %%%
\noindent\textbf{Question 1: Identify the pieces.}  Write down what the running variable $X_i$, the cutoff $c$, the treatment $D_i$, and a candidate outcome $Y_i$ are in this design.  Then state, in one sentence, what the continuity assumption requires here.

\vspace{6cm}

%%% QUESTION 2 %%%
\noindent\textbf{Question 2: What else changes at the cutoff?}  The biggest threat to RD is another treatment at the same threshold (``bundled treatments'').  Name \emph{at least one} other thing that changes at age 65 that could plausibly affect hospital utilization, and explain in a sentence why it is a problem for the RD interpretation.

\vspace{6cm}

\newpage

%%% QUESTION 3 %%%
\noindent\textbf{Question 3: Reading a balance plot.}  Suppose you run an RD on Brazilian municipalities at the e-voting threshold (40,500 registered voters) and produce balance plots like the one we saw in lecture.  For each finding below, say in one sentence whether it \emph{supports} the design, \emph{threatens} the design, or is \emph{ambiguous}, and why.

\medskip
\begin{enumerate}[label=(\alph*), itemsep=0.5em]
\item Pre-1998 municipal income shows no jump at the cutoff.

\vspace{2.2cm}

\item Pre-1998 literacy rate shows a small but statistically significant jump at the cutoff.

\vspace{2.2cm}

\item \emph{Post-1998} federal social spending shows a jump at the cutoff.

\vspace{2.2cm}
\end{enumerate}

%%% QUESTION 4 %%%
\noindent\textbf{Question 4: Sharp vs.\ fuzzy.}  For each scenario below, say whether you'd treat it as a \textbf{sharp} or \textbf{fuzzy} RD, and -- if fuzzy -- what the estimand becomes (and for whom).

\medskip
\begin{enumerate}[label=(\alph*), itemsep=0.5em]
\item French municipal elections: cities with population $\geq 3{,}500$ use proportional representation; below $3{,}500$ they use first-past-the-post.  Outcome: number of effective parties.

\vspace{2.2cm}

\item A scholarship program grants automatic offers to students scoring above an SAT threshold; students just below the threshold can still be considered, but the probability of an offer drops sharply.  Outcome: college completion.

\vspace{2.2cm}
\end{enumerate}

\end{document}
